\documentclass[11pt]{article}
\usepackage{amsmath,amsthm,amssymb}
\usepackage[margin=1in]{geometry}
\usepackage{enumitem}
\usepackage{hyperref}

\newtheorem{theorem}{Theorem}
\newtheorem{lemma}[theorem]{Lemma}
\newtheorem{proposition}[theorem]{Proposition}
\newtheorem{corollary}[theorem]{Corollary}
\theoremstyle{definition}
\newtheorem{definition}[theorem]{Definition}
\newtheorem{remark}[theorem]{Remark}

\title{The Shift Lemma and the $j$-formula for all hooks:\\
       structural proof by induction on the arm length}
\author{Clio}
\date{13 May 2026 (very late evening)}

\begin{document}
\maketitle

\begin{abstract}
We isolate a clean operator identity inside the Iwahori--Hecke algebra
$H_q(S_n)$ --- the \emph{Shift Lemma}
\[
   R'_a R'_{a+1} \cdots R'_b \;=\; (T_{a-1}{+}1)(T_a{+}1)\cdots(T_{b-1}{+}1)\,
   R'_{a-1} R'_a \cdots R'_{b-1}, \qquad 2 \le a \le b \le n
\]
--- and use it to give a structural proof of the $j$-formula for
\textbf{every hook} $\lambda = (k, 1^{n-k})$ in the rank-zero regime
$n \ge 2k$.

\textbf{Theorem (main).} For every $k \ge 2$ and every $n \ge 2k$,
\[
   B^{(\lambda)}_{j_0(\lambda)}\,V_\lambda
   \;:=\; R'_{j_0(\lambda)}\,R'_{j_0(\lambda)+1}\,\cdots\,R'_n\,V_\lambda
   \;\subseteq\; E_1^-\!\mid_{V_\lambda},
   \qquad
   j_0(\lambda) = n - k + 2.
\]
In particular $\Pi^{S_n}|_{V_\lambda} = 0$ unconditionally for every hook
$\lambda$ in the rank-zero regime.

\textbf{The proof is a clean induction on $k$.}  Base cases $k = 2$
(May-11) and $k = 1$ (sign rep, trivial).  Inductive step:
the Phase~A decomposition $E_{n-1}^+|_{V_\lambda} = \Phi_*(V_{\mu_1}) \oplus S$
splits the image $R'_n V_\lambda$ into two pieces:
\begin{itemize}[leftmargin=2em,topsep=0.2em,itemsep=0.1em]
\item The $\Phi_*(V_{\mu_1})$ piece, with $\mu_1 = (k-1, 1^{n-k-1})$, is
handled by the May-13 r1 commutation (which carries through hooks
unconditionally) combined with the inductive hypothesis at $k - 1$.
\item The same-row part $S$, isomorphic as an $H_q(S_{n-2})$-module to
$V_{\nu'}$ with $\nu' = (k-2, 1^{n-k})$, is killed by an application of the
Shift Lemma: the iteration $R'_{j_0}\cdots R'_{n-1}|_S$ rearranges into
$T$-factors $\times$ an inner iteration that, via the iso $S \cong V_{\nu'}$
and the inductive hypothesis at $k - 2$, equals zero.
\end{itemize}

The Shift Lemma is the operator-theoretic engine: it lets us \emph{shift
the entire iteration down by one in indices}, at the cost of multiplying
through by a row of commuting $T$-factors.  This shift-by-one is exactly
what's needed to align the same-row iteration on $\lambda$ with the sharp
$j$-formula iteration on $\nu'$.

\textbf{Consequences.}  Beyond the hook $j$-formula itself:
\begin{enumerate}[leftmargin=2em,topsep=0.2em,itemsep=0.1em]
\item The May-12 evening-3 paper's $\dim B = 3$ for
$\lambda = (4, 2, 1^{n-6})$ becomes \emph{unconditional} (the upper bound
relied on the hook $j$-formula at $k = 4$, which was conditional then;
this paper closes that gap).
\item The same-row-dies conjecture for $\lambda = (k, k, 1^q)$ reduces
to the $j$-formula for the sub-shape $\nu = (k, k-2, 1^q)$ via the same
Shift Lemma machinery (concrete next target on this thread).
\item The framework is shape-agnostic: any same-row-dies argument can
be carried out via Shift Lemma plus a $j$-formula on the sub-shape.
\end{enumerate}

The mathematical content is short and self-contained; the heaviest input
is the May-11/May-12 hook base cases and the Phase~A endpoint (May-19).
\end{abstract}

\section{Setup}\label{sec:setup}

Notation follows the May-13 r1 paper
(\texttt{2026-05-13-r1-inductive-reduction.tex}) and the May-12 hook
paper (\texttt{2026-05-12-hook-j-formula.tex}).
$H_q(S_n)$ is the Iwahori--Hecke algebra over $\mathbb{Q}(q)$ with
generators $T_1, \ldots, T_{n-1}$ satisfying the usual quadratic
relations $(T_i - q)(T_i + 1) = 0$ and the braid relations.
For $\lambda \vdash n$, $V_\lambda$ is the seminormal cell module with
basis $\{v_T : T \in \mathrm{SYT}(\lambda)\}$.

\begin{definition}\label{def:notation}
For $1 \le k \le n$, set
\[
   R'_k \;:=\; (T_{k-1}{+}1)\,(T_{k-2}{+}1)\,\cdots\,(T_1{+}1)
   \;\in\; H_q(S_n),
\]
with the convention $R'_1 = 1$ (empty product).  For $1 \le a \le b \le n$,
set
\[
   P_{a, b} \;:=\; R'_a\,R'_{a+1}\,\cdots\,R'_b,
\]
with the convention $P_{a, b} = 1$ when $a > b$.
For a partition $\lambda$,
\[
   B^{(\lambda)}_j\,V_\lambda \;:=\; P_{j, n}\,V_\lambda
   \;=\; R'_j\,R'_{j+1}\,\cdots\,R'_n\,V_\lambda,
   \qquad
   j_0(\lambda) := \ell(\lambda) + 1.
\]
Write $E_j^\pm$ for the $\pm$-eigenspaces of $T_j$ on $V_\lambda$
(eigenvalues $q$ resp.\ $-1$).
\end{definition}

\section{The Shift Lemma}\label{sec:shift}

\begin{lemma}[Shift Lemma]\label{lem:shift}
For $2 \le a \le b \le n$, inside $H_q(S_n)$,
\begin{equation}\label{eq:shift}
   P_{a, b}
   \;=\;
   (T_{a-1}+1)\,(T_a+1)\,\cdots\,(T_{b-1}+1)\;\cdot\; P_{a-1, b-1}.
\end{equation}
\end{lemma}

\begin{proof}
By induction on $b - a$.

\medskip\noindent\textbf{Base case $b = a$:}  $P_{a, a} = R'_a =
(T_{a-1}+1)\,R'_{a-1} = (T_{a-1}+1)\,P_{a-1, a-1}$.  Matches~\eqref{eq:shift}.

\medskip\noindent\textbf{Inductive step:}  Assume~\eqref{eq:shift} for
$P_{a, b-1}$ (i.e., the lemma for length $b - 1 - a$).  We have
\[
   P_{a, b} \;=\; P_{a, b-1}\,R'_b
   \;\stackrel{\text{IH}}{=}\;
   (T_{a-1}+1)\,(T_a+1)\,\cdots\,(T_{b-2}+1)\,P_{a-1, b-2}\,R'_b.
\]
Unfold $R'_b = (T_{b-1}+1)\,R'_{b-1}$:
\begin{equation}\label{eq:step}
   P_{a, b}
   \;=\;
   (T_{a-1}+1)\,(T_a+1)\,\cdots\,(T_{b-2}+1)\,P_{a-1, b-2}\,(T_{b-1}+1)\,R'_{b-1}.
\end{equation}

\textbf{Key commutation.} $P_{a-1, b-2}$ commutes with $(T_{b-1}+1)$.
Indeed, $P_{a-1, b-2} = R'_{a-1}\,R'_a\,\cdots\,R'_{b-2}$ is a product of
factors $(T_j{+}1)$ for $j \le b - 3$ (the maximum index of a $T$-factor
inside $R'_{b-2}$ is $b - 3$).  Each such $T_j$ has $|j - (b-1)| \ge 2$,
so $T_j$ commutes with $T_{b-1}$, hence so does the product.
Therefore~\eqref{eq:step} becomes
\[
   P_{a, b}
   \;=\;
   (T_{a-1}+1)\,(T_a+1)\,\cdots\,(T_{b-2}+1)\,(T_{b-1}+1)\,P_{a-1, b-2}\,R'_{b-1}.
\]
Since $P_{a-1, b-2}\,R'_{b-1} = P_{a-1, b-1}$ by definition, we
get~\eqref{eq:shift} for $P_{a, b}$.
\end{proof}

\begin{remark}\label{rem:shift-special-case}
Specialising to $b = a + 1$ recovers the commutation identity used in
the May-13 night same-row-dies paper for $(3, 3, 1^q)$:
\[
   R'_a\,R'_{a+1}
   \;=\;
   (T_{a-1}+1)\,(T_a+1)\,R'_{a-1}\,R'_a.
\]
The Shift Lemma is the natural generalisation to chains of arbitrary
length.
\end{remark}

\begin{remark}\label{rem:shift-computational}
The Shift Lemma~\eqref{eq:shift} holds as a literal operator identity
inside $H_q(S_n)$ (independent of any representation), verified
computationally at $P = 100\,003$, $Q = 11$ in:
\texttt{\textasciitilde/projects/scratch/2026-05-13-shift-lemma/verify.py}
for $\lambda \in \{(3, 2), (2, 2, 1), (3, 1^3), (4, 1^3), (2, 2, 1^2),
(3, 2, 1^2)\}$ at $n \in \{5, 6, 7\}$ and all valid $(a, b)$ pairs.
\end{remark}

\section{Hook structure and prior ingredients}\label{sec:ingredients}

Fix $\lambda = (k, 1^{n-k})$ with $k \ge 2$ and $n \ge 2k$.  We have
$\ell(\lambda) = n - k + 1$ and $j_0(\lambda) = n - k + 2$.  The
removable corners of $\lambda$ are:
\[
   c_1 := (1, k) \quad\text{(length-preserving)},
   \qquad
   c_* := (n - k + 1, 1) \quad\text{(column-1 bottom)}.
\]

\subsection{Phase A and the $E_+$ decomposition}\label{sec:phaseA}

\begin{theorem}[Phase A endpoint, May-19]\label{thm:phaseA}
For every $H_q(S_n)$-module $V$ and every $1 \le j \le n - 1$,
$(T_j+1)\,(T_{j-1}+1)\,\cdots\,(T_1+1)\,V = E_j^+(V)$.  In particular,
$R'_n\,V_\lambda = E_{n-1}^+\!\mid_{V_\lambda}$.
\end{theorem}

\begin{lemma}[Decomposition of $E_+$ for hooks]\label{lem:E-decomp}
For $\lambda = (k, 1^{n-k})$ with $k \ge 3$,
\begin{equation}\label{eq:E-decomp}
   E_{n-1}^+\!\mid_{V_\lambda}
   \;=\;
   \Phi_*(V_{\mu_1}) \;\oplus\; S,
\end{equation}
where
\begin{itemize}[leftmargin=2em,topsep=0.2em,itemsep=0.1em]
\item $\mu_1 := \lambda \setminus \{c_1, c_*\} = (k-1, 1^{n-k-1})$
on $n - 2$ letters;
\item $\Phi_*: V_{\mu_1} \hookrightarrow V_\lambda$ is the
$+q$-eigenvector embedding of the $T_{n-1}$-$2$-block at the corner
pair $\{c_1, c_*\}$ (May-13 r1 Defn~2.1);
\item $S$ is the span of seminormal basis vectors $v_T$ for SYTs of
$\lambda$ with the letters $\{n-1, n\}$ at the same-row pair
$\{(1, k-1), (1, k)\}$ at row $1$.
\end{itemize}
For $k = 2$, $S = 0$ and $E_+ = \Phi_*(V_{\mu_1})$ already.
\end{lemma}

\begin{proof}
$T_{n-1}$ acts on $v_T$ by the Hoefsmit rule depending on the relative
positions of $n - 1, n$ in $T$.  Eigenvalue $+q$ arises in two cases:
\begin{enumerate}[leftmargin=2em,topsep=0.2em,itemsep=0.1em,label=(\arabic*)]
\item $\{n-1, n\}$ at a corner pair (one length-pres, one column-1 bottom)
in different row and column: $T_{n-1}$ is a $2 \times 2$ Hoefsmit block,
one-dim $+q$-eigenvector inside.  The only such pair for a hook is
$\{c_1, c_*\}$.  The $+q$-eigenvectors are exactly the $\Phi_*(v_{T'})$
for $T' \in \mathrm{SYT}(\mu_1)$ (May-13 r1 Defn~2.1).
\item $\{n-1, n\}$ at a same-row pair $\{(i, \lambda_i - 1), (i, \lambda_i)\}$:
$T_{n-1} v_T = q v_T$.
\end{enumerate}
For a hook $\lambda = (k, 1^{n-k})$, the same-row pair is valid at row
$i$ iff $\lambda_i \ge 2$ and ($\lambda_{i+1} \le \lambda_i - 2$ or
$i = \ell$).  Only row $1$ satisfies $\lambda_1 = k \ge 2$, and the
condition $\lambda_2 = 1 \le k - 2$ holds iff $k \ge 3$.  Hence
for $k \ge 3$ there is a unique same-row pair at row $1$,
$\{(1, k-1), (1, k)\}$, and the corresponding $S$ is the span of
seminormal vectors with $n-1, n$ at this pair.  The remaining cells of
$\lambda$ form the sub-partition
$\nu' := \lambda \setminus \{(1, k-1), (1, k)\} = (k-2, 1^{n-k})$ on
$n - 2$ letters.

The direct sum~\eqref{eq:E-decomp} is by construction (distinct
seminormal-basis supports), and orthogonal to the seminormal pairing.
\end{proof}

\subsection{The $\Phi_*$-commutation}\label{sec:Phi-commute}

\begin{lemma}[$\Phi_*$ commutation, May-13 r1 Lemma~3.1]\label{lem:Phi-commute}
For any partition $\lambda$ with a single length-pres corner $c_1$ and
column-1 bottom $c_*$ (hypotheses (H1)--(H3) of the May-13 r1 paper),
\[
   R'_{\ell+1}\,R'_{\ell+2}\,\cdots\,R'_{n-1}\,\big|_{\Phi_*(V_{\mu_1})}
   \;=\;
   (T_\ell+1)\,(T_{\ell+1}+1)\,\cdots\,(T_{n-2}+1)\,
   \Phi_*\big(R'^{(\mu_1)}_\ell\,R'^{(\mu_1)}_{\ell+1}\,\cdots\,R'^{(\mu_1)}_{n-2}\big).
\]
In particular,
\[
   R'_{\ell+1}\,\cdots\,R'_{n-1}\,\Phi_*(V_{\mu_1})
   \;=\;
   (T_\ell+1)\,\cdots\,(T_{n-2}+1)\,\Phi_*\big(B^{(\mu_1)}_{j_0(\mu_1)}\,V_{\mu_1}\big).
\]
\end{lemma}

\begin{proof}
See May-13 r1 paper, Section~3 and the May-15 paper.  The proof is
purely operator-theoretic: $\Phi_*$ is $H_q(S_{n-2})$-equivariant under
$\langle T_1, \ldots, T_{n-3}\rangle$, and the iteration uses only
$(T_j+1)$ factors with $j \le n - 2$, all of which (after the initial
$(T_{n-2}+1)$) commute through $\Phi_*$ via index-gap arguments.
Hooks satisfy (H1)--(H3): one length-pres corner $c_1 = (1, k)$,
column-1 bottom $c_* = (n-k+1, 1)$, axial distance
$d = -(n - k) - (k - 1) = -(n - 1)$, $|d| = n - 1 \ge 2$ for $n \ge 3$.
\end{proof}

\subsection{The same-row iso $\Psi$}\label{sec:Psi}

\begin{lemma}[Same-row iso]\label{lem:Psi-iso}
For $\lambda = (k, 1^{n-k})$ with $k \ge 3$, the map
\[
   \Psi: S \;\xrightarrow{\sim}\; V_{\nu'}, \qquad \nu' = (k - 2, 1^{n-k})
   \text{ on } n - 2 \text{ letters},
\]
sending $v_T \in S$ (where $T$ has $n-1, n$ at $(1, k-1), (1, k)$) to
$v_{T|_{\nu'}}$ (the restriction of $T$ to the letters $\{1, \ldots, n-2\}$
filling the cells of $\nu'$) is a $\mathbb{Q}(q)$-linear isomorphism,
$H_q(S_{n-2})$-equivariant under $\langle T_1, \ldots, T_{n-3}\rangle$.
\end{lemma}

\begin{proof}
For $T \in S$, the letters $n-1, n$ are at $(1, k-1), (1, k)$, which are
NOT cells of $\nu'$.  The letters $\{1, \ldots, n-2\}$ fill the cells
of $\nu'$ (namely row $1$ of length $k - 2$ and the column-$1$ tail of
length $n - k + 1$).  So $T|_{\nu'}$ is a well-defined SYT of $\nu'$.
The map is bijective by construction.

For equivariance: for $j \le n - 3$, $T_j$ acts on letters $j, j+1$,
both $\le n - 2$, hence in cells of $\nu'$.  The Hoefsmit action of
$T_j$ on $v_T$ depends only on the positions of $j, j+1$ in $T$, which
coincide with their positions in $T|_{\nu'}$.  Hence
$T_j \Psi^{-1}(v) = \Psi^{-1}(T_j v)$.
\end{proof}

\begin{remark}\label{rem:E1-preservation}
Since $\Psi$ is $T_1$-equivariant (because $1 \le n - 3$ for $n \ge 4$),
$\Psi$ preserves the $T_1$-eigenspaces.  In particular $\Psi(S \cap E_1^-)
= E_1^-|_{V_{\nu'}}$.  Similarly for $\Phi_*$ (May-13 r1 Lemma~2.1):
$\Phi_*$ is $T_1$-equivariant for $n \ge 4$, so $\Phi_*(E_1^-|_{V_{\mu_1}})
\subseteq E_1^-|_{V_\lambda}$.
\end{remark}

\section{Main theorem and proof}\label{sec:main}

\begin{theorem}[Hook $j$-formula, all $k$]\label{thm:main}
For every $k \ge 1$ and every $n \ge 2k$, the hook $\lambda = (k, 1^{n-k})$
satisfies
\begin{equation}\label{eq:main}
   B^{(\lambda)}_{j_0(\lambda)}\,V_\lambda
   \;\subseteq\; E_1^-\!\mid_{V_\lambda},
   \qquad
   j_0(\lambda) = n - k + 2.
\end{equation}
\end{theorem}

\begin{proof}
Strong induction on $k$.

\medskip\noindent\textbf{Base case $k = 1$:}
$\lambda = (1^n)$ is the sign representation; $T_i$ acts as $-1$ for
every $i$.  So $R'_j V_\lambda = 0$ for every $j \ge 2$ (each factor
$(T_i+1) = 0$ on $V_\lambda$).  Hence $B^{(\lambda)}_j V_\lambda = 0$
trivially $\subseteq E_1^-|_{V_\lambda}$, for every $j$.

\medskip\noindent\textbf{Base case $k = 2$:}
$\lambda = (2, 1^{n-2})$, May-11 paper
(\texttt{2026-05-11-j-formula.tex}, Theorem on $\lambda = (2, 1^{n-2})$).

\medskip\noindent\textbf{Inductive step $k \ge 3$:}
Assume~\eqref{eq:main} for all $k' < k$ (in their rank-zero regime).
Let $\mu_1 = (k-1, 1^{n-k-1})$ and $\nu' = (k-2, 1^{n-k})$, both on
$n - 2$ letters.  Both are hooks; rank-zero regime for $\mu_1$ is
$n - 2 \ge 2(k-1)$, i.e., $n \ge 2k$ (satisfied); similarly $\nu'$ at
$n - 2 \ge 2(k-2) = 2k - 4$, $n \ge 2k - 2$ (weaker; satisfied).
By IH:
\begin{equation}\label{eq:IH}
   B^{(\mu_1)}_{j_0(\mu_1)}\,V_{\mu_1} \;\subseteq\; E_1^-\!\mid_{V_{\mu_1}},
   \qquad
   B^{(\nu')}_{j_0(\nu')}\,V_{\nu'} \;\subseteq\; E_1^-\!\mid_{V_{\nu'}}.
\end{equation}

By Phase A (Theorem~\ref{thm:phaseA}) and the $E_+$ decomposition
(Lemma~\ref{lem:E-decomp}),
\[
   R'_n\,V_\lambda \;=\; E_{n-1}^+\!\mid_{V_\lambda}
   \;=\; \Phi_*(V_{\mu_1}) \;\oplus\; S.
\]
Hence
\[
   B^{(\lambda)}_{j_0(\lambda)}\,V_\lambda
   \;=\; R'_{j_0}\,R'_{j_0+1}\,\cdots\,R'_{n-1}\,R'_n\,V_\lambda
   \;=\; R'_{j_0}\,\cdots\,R'_{n-1}\,\big(\Phi_*(V_{\mu_1}) + S\big).
\]
We show each summand is in $E_1^-|_{V_\lambda}$.

\medskip\noindent\textbf{Part 1: $\Phi_*$-piece is in $E_1^-$.}

By Lemma~\ref{lem:Phi-commute} applied to $\lambda$ (hook, $r = 1$):
\[
   R'_{j_0}\,\cdots\,R'_{n-1}\,\Phi_*(V_{\mu_1})
   \;=\;
   (T_\ell+1)\,(T_{\ell+1}+1)\,\cdots\,(T_{n-2}+1)\,
   \Phi_*\!\left(B^{(\mu_1)}_{j_0(\mu_1)}\,V_{\mu_1}\right),
\]
where $\ell = n - k + 1$ and $j_0(\mu_1) = (n-2) - (k-1) + 2 = n - k + 1$.
By the IH on $\mu_1$ (first part of~\eqref{eq:IH}),
$B^{(\mu_1)}_{j_0(\mu_1)} V_{\mu_1} \subseteq E_1^-|_{V_{\mu_1}}$.  By
Remark~\ref{rem:E1-preservation}, $\Phi_*$ preserves $E_1^-$, so
$\Phi_*(B^{(\mu_1)}_{j_0(\mu_1)} V_{\mu_1}) \subseteq E_1^-|_{V_\lambda}$.
The outer $T$-factors $(T_\ell+1), (T_{\ell+1}+1), \ldots, (T_{n-2}+1)$
all have index $\ge \ell = n - k + 1 \ge 3$ (since $n \ge 2k \ge k + 2$
for $k \ge 2$), hence commute with $T_1$ (index gap $\ge 2$) and
preserve $E_1^-|_{V_\lambda}$.  Therefore
\[
   R'_{j_0}\,\cdots\,R'_{n-1}\,\Phi_*(V_{\mu_1})
   \;\subseteq\; E_1^-\!\mid_{V_\lambda}.
\]

\medskip\noindent\textbf{Part 2: $S$-piece dies (same-row dies).}

By the Shift Lemma (Lemma~\ref{lem:shift}) with $a = j_0$ and $b = n - 1$,
\begin{equation}\label{eq:shift-app}
   R'_{j_0}\,\cdots\,R'_{n-1}
   \;=\;
   (T_{j_0 - 1}+1)\,(T_{j_0}+1)\,\cdots\,(T_{n-2}+1)\,
   R'_{j_0 - 1}\,R'_{j_0}\,\cdots\,R'_{n-2}.
\end{equation}
Note $j_0 - 1 = n - k + 1 = \ell$, and the right-hand inner factor is
$P_{\ell, n-2}$.

Apply to $S$.  The operators $R'_j$ for $j \le n - 2$ all involve only
$T_i$ with $i \le n - 3$, so by Lemma~\ref{lem:Psi-iso}
(equivariance of $\Psi$ under $\langle T_1, \ldots, T_{n-3}\rangle$),
\[
   P_{\ell, n-2}\,\big|_S
   \;=\; \Psi^{-1}\big(\,P^{(\nu')}_{\ell, n-2}\,V_{\nu'}\,\big)
   \;=\; \Psi^{-1}\big(\,B^{(\nu')}_\ell\,V_{\nu'}\,\big),
\]
where the rightmost identification is just unwinding the notation: for
$\nu' = (k-2, 1^{n-k})$ on $n - 2$ letters,
\[
   \ell(\nu') = 1 + (n - k) = n - k + 1 = \ell,
   \qquad
   j_0(\nu') = n - k + 2 = j_0(\lambda),
\]
and $|\nu'| = n - 2$, so
$P^{(\nu')}_{\ell, n-2} = R'^{(\nu')}_\ell R'^{(\nu')}_{\ell+1} \cdots R'^{(\nu')}_{|\nu'|}
= B^{(\nu')}_{\ell(\nu')}\,V_{\nu'}$.

This is the $B$-image \emph{one step below} the sharp $j_0(\nu')$:
\begin{equation}\label{eq:nu-killed}
   B^{(\nu')}_{\ell(\nu')}\,V_{\nu'}
   \;=\;
   R'^{(\nu')}_{\ell(\nu')}\,\cdot\,B^{(\nu')}_{j_0(\nu')}\,V_{\nu'}
   \;\subseteq\;
   R'^{(\nu')}_{\ell(\nu')}\,\cdot\,E_1^-\!\mid_{V_{\nu'}}
   \;=\; 0.
\end{equation}
The middle inclusion is the IH on $\nu'$ (second part of~\eqref{eq:IH});
the final equality is because $R'^{(\nu')}_{\ell(\nu')}$ has rightmost
factor $(T_1+1)$, which annihilates $E_1^-|_{V_{\nu'}}$.

Therefore $P_{\ell, n-2}|_S = \Psi^{-1}(0) = 0$, so by~\eqref{eq:shift-app},
\[
   R'_{j_0}\,\cdots\,R'_{n-1}\,S
   \;=\;
   (T_{j_0-1}+1)\,\cdots\,(T_{n-2}+1)\,P_{\ell, n-2}\,S
   \;=\; 0.
\]

\medskip\noindent\textbf{Conclusion:}  Combining Parts 1 and 2,
$B^{(\lambda)}_{j_0(\lambda)} V_\lambda \subseteq E_1^-|_{V_\lambda}
+ 0 = E_1^-|_{V_\lambda}$.  This closes the induction.
\end{proof}

\section{Consequences}\label{sec:consequences}

\begin{corollary}[Rank-zero theorem for all hooks]\label{cor:rank-zero}
For every $k \ge 2$ and every $n \ge 2k$,
\[
   \Pi^{S_n}\!\mid_{V_{(k, 1^{n-k})}} \;=\; 0.
\]
\end{corollary}

\begin{proof}
By Theorem~\ref{thm:main}, $B^{(\lambda)}_{j_0} V_\lambda \subseteq E_1^-$.
Iteratively descending: $B^{(\lambda)}_{j_0 - 1} = R'_{j_0 - 1}\,
B^{(\lambda)}_{j_0} \subseteq R'_{j_0 - 1}\,E_1^- = 0$ (rightmost
$(T_1+1)$ kills $E_1^-$).  So $B^{(\lambda)}_j = 0$ for every
$2 \le j \le j_0 - 1$; in particular $\Pi^{S_n}|_{V_\lambda} =
B^{(\lambda)}_2 = 0$.
\end{proof}

\begin{corollary}[$\dim B^{(\lambda)}_{j_0} \le 1$ for all hooks]\label{cor:dim-bound}
For every $k \ge 2$ and every $n \ge 2k$,
\[
   \dim B^{(\lambda)}_{j_0(\lambda)}\,V_\lambda \;\le\; 1.
\]
\end{corollary}

\begin{proof}
Same-row dies (Part 2 of the proof of Theorem~\ref{thm:main}) gives
the equality (not just inclusion)
\[
   B^{(\lambda)}_{j_0(\lambda)}\,V_\lambda
   \;=\;
   (T_\ell+1)\,\cdots\,(T_{n-2}+1)\,\Phi_*\!\big(B^{(\mu_1)}_{j_0(\mu_1)}\,V_{\mu_1}\big)
\]
(May-13 r1 Theorem~3, unconditional now that same-row-dies is proved
for hooks).  Since $\Phi_*$ is injective and outer $T$-factors don't
increase dimension, $\dim B^{(\lambda)}_{j_0} \le \dim B^{(\mu_1)}_{j_0(\mu_1)}$.
By induction down to $k = 2$ where $\dim B^{(2, 1^{m-2})}_{j_0} = 1$
(May-11), the bound $\le 1$ propagates.
\end{proof}

\begin{remark}[On $\dim = 1$]\label{rem:dim-1}
The matching lower bound $\dim B^{(\lambda)}_{j_0} \ge 1$ (giving
$\dim = 1$ and the Catalan-many support of the May-12 hook paper's
Conjecture~5.1) is a separate question --- it requires showing the
outer $T$-factors are not annihilating on the $\Phi_*$-image.  Direct
seminormal-basis tracking establishes $\dim = 1$ for $k \le 4$ at
$n \le 9$ (May-12, May-12 evening-3); the full structural proof at
general $(k, n)$ is unaddressed here.  This paper focuses on the
$E_1^-$-containment, which is the $j$-formula proper.
\end{remark}

\begin{corollary}[Unconditional $j$-formula on $(4, 2, 1^{n-6})$]\label{cor:421}
The May-12 evening-3 paper's $j$-formula for $\lambda = (4, 2, 1^{n-6})$
at $n \ge 10$ is now \emph{unconditional}:
$B^{(\lambda)}_{n-3}\,V_\lambda \subseteq E_1^-\!\mid_{V_\lambda}$,
and $\dim B^{(\lambda)}_{n-3}\,V_\lambda = 3$.
\end{corollary}

\begin{proof}
The May-12 evening-3 paper assumes (Theorem~3.2 of that paper) the
hook $j$-formula at $k = 4$ for $\mu_B = (4, 1^{n-6})$ on $n - 2$
letters.  By Theorem~\ref{thm:main} (at $k = 4$), this is now
established structurally.  The rest of the May-12 evening-3 argument
goes through verbatim, giving the upper bound $\le 3$ and (with the
explicit construction) the matching lower bound.
\end{proof}

\section{Computational verification}\label{sec:verification}

\paragraph{Scripts.}
\texttt{\textasciitilde/projects/scratch/2026-05-13-shift-lemma/verify.py}
(Shift Lemma + hook $j$-formula) and
\texttt{verify\_extras.py} (same-row-dies on $(k, k, 1^q)$, larger
hooks).  Mod $P = 100\,003$, $Q = 11$.

\paragraph{Shift Lemma.}  Verified as a literal operator identity in
$\mathrm{End}(V_\lambda)$ for $\lambda \in \{(3, 2), (2, 2, 1),
(3, 1^3), (4, 1^3), (2, 2, 1^2), (3, 2, 1^2)\}$ at all valid $(a, b)$
pairs in $2 \le a \le b \le n$.  All pass.

\paragraph{Hook $j$-formula.}  Verified for $(k, n) \in
\{(2, 4), (2, 5), (2, 6), (3, 6), (3, 7), (3, 8), (4, 8), (4, 9),
(4, 10), (5, 10)\}$.  In every case, $(T_1+1)$ applied to
$B^{(\lambda)}_{j_0} V_\lambda$ has rank $0$, confirming the $E_1^-$
containment.

\paragraph{Same-row dies for $(k, k, 1^q)$, sanity.}  Same-row part
$S$ has the predicted dimension $f^{(k, k-2, 1^q)}$ and the prediction
$P_{j_0(\lambda), n-1}\,S = 0$ holds for $(3, 3, 1^q)$ at
$q \in \{2, 3, 4, 5\}$ (recovering the May-13 night theorem).  For
$(4, 4, 1^q)$, the prediction holds only when $q \ge 4$ (the regime
where $\nu = (4, 2, 1^q)$ enters its own rank-zero regime); at $q = 2, 3$,
$\mathrm{rank}(P_{j_0, n-1}\,S) = 3$, consistent with the failure of
$\nu$'s $j$-formula in its boundary regime.  This sharpens what
"same-row dies" actually requires: the $j$-formula on $\nu$, which is
itself shape- and regime-sensitive.

\section{Beyond hooks: the framework's reach}\label{sec:beyond}

The proof of Theorem~\ref{thm:main} factors into three independent
ingredients:
\begin{enumerate}[leftmargin=2em,topsep=0.2em,itemsep=0.1em,label=(\arabic*)]
\item \textbf{Phase A + $E_+$ decomposition}: $R'_n V_\lambda = E_+$ as
a direct sum of $\Phi_X$-pieces (multi-corner) and $S_i$-pieces
(same-row).  General fact.
\item \textbf{$\Phi$-commutation}: each multi-corner piece is handled by
the May-13 r1 / May-15 $\Phi$-commutation (combined with the
multi-corner reduction of May-12 evening-3).
\item \textbf{Shift Lemma + iso}: each same-row piece reduces to the
$j$-formula on the sub-shape $\nu_i' = \lambda \setminus \{\text{row-}i
\text{ same-row pair}\}$ via the Shift Lemma and the iso $\Psi_i: S_i
\xrightarrow{\sim} V_{\nu_i'}$.
\end{enumerate}

The induction on $k$ for hooks works because both subproblems
($\mu_1$ and $\nu'$) are hooks with smaller arm length.  For general
shapes, the induction descends along multi-corner sub-shapes and
same-row sub-shapes; both decrease the size of the "decoration shape"
$\lambda^{(\ge 2)}$ (the column-$1$-deleted partition).

\paragraph{Concrete next applications.}
\begin{itemize}[leftmargin=2em,topsep=0.2em,itemsep=0.1em]
\item \textbf{$j$-formula on $(k, 2, 1^q)$ for $k \ge 4$.}  $r = 2$
multi-corner shape with $\mu_A = (k-1, 2, 1^{q-1})$ (recursive),
$\mu_B = (k, 1^{q-1})$ (hook, this paper).  The May-12 evening-3
machinery handles $k = 4$; iterating in $k$ should work.
\item \textbf{Same-row dies on $(k, k, 1^q)$ for $k = 4$, $q \ge 4$.}
Reduces via the Shift Lemma to the $j$-formula on
$\nu = (k, k-2, 1^q)$ in $\nu$'s rank-zero regime.  For $k = 4$, $\nu =
(4, 2, 1^q)$ on $m = q + 6$ letters; rank-zero requires
$q \ge 4$.  In that regime, Corollary~\ref{cor:421} supplies the
needed $j$-formula on $\nu$, so same-row dies for $(4, 4, 1^q)$ at
$q \ge 4$ follows.  Below $q = 4$, $\nu$ is at or below the rank-zero
boundary and the same-row part does \emph{not} die (verified
computationally at $(4, 4, 1^2)$ and $(4, 4, 1^3)$,
where $\mathrm{rank}(P_{j_0, n-1}\,S) = 3 \ne 0$).
\item \textbf{$j$-formula on $(k, k-2, 1^q)$ for $k \ge 5$.}  Same-row
pair on row $2$ gives $\nu' = (k, k-4, 1^q)$; recursion in $k - 2$.
Requires also the $\Phi$-piece argument with $\mu_1$, $\mu_2$
themselves $r = 2$ non-hook shapes.
\end{itemize}

The hook case (this paper) is the cleanest application of the
framework; the general theorem~``$j$-formula in maximum generality, in
the stable regime'' is within reach but requires combining all three
ingredients carefully across nested inductions.

\section{Honest status}\label{sec:status}

\paragraph{What's proved unconditionally.}  Theorem~\ref{thm:main}: the
$j$-formula $B^{(\lambda)}_{j_0(\lambda)} V_\lambda \subseteq
E_1^-\!\mid_{V_\lambda}$ for every hook $\lambda = (k, 1^{n-k})$ with
$k \ge 1$ and $n \ge 2k$.  This subsumes May-11 ($k = 2$) and May-12
($k = 3$) and extends to all $k$.

\paragraph{What's a corollary.}  Rank-zero theorem for all hooks in the
rank-zero regime (Corollary~\ref{cor:rank-zero}); $\dim
B^{(\lambda)}_{j_0} \le 1$ (Corollary~\ref{cor:dim-bound});
unconditional $j$-formula on $(4, 2, 1^{n-6})$
(Corollary~\ref{cor:421}).

\paragraph{What's not proved.}  The matching lower bound $\dim = 1$
(May-12 paper's Conjecture~5.1, the Catalan support phenomenon).  My
framework gives $\dim \le 1$ via the same-row-dies machinery, but the
non-vanishing of the outer $T$-factor action on the $\Phi_*$-image is
not established here.

\paragraph{What it took.}  The proof is short --- three pages of
ingredients, two pages of induction.  The key was identifying the
Shift Lemma (Section~\ref{sec:shift}) as the operator-theoretic
mechanism that aligns the same-row iteration on $\lambda$ with the
sharp $j$-formula iteration on $\nu'$.  Once the Shift Lemma is in
hand, the rest is routine induction.

The Shift Lemma is also the right framework for the more general
same-row-dies thread: any shape $\lambda$ with a same-row pair admits
the same one-step reduction, modulo a $j$-formula on the sub-shape.

\end{document}
